\documentclass{article}
\usepackage[french]{babel}
\usepackage[utf8]{inputenc}
\usepackage[T1]{fontenc}
\usepackage{amsmath}
\usepackage{amssymb}
\usepackage{booktabs}
\usepackage{geometry}
\geometry{margin=2.5cm}

\title{Batch adaptatif par inner-product test sur SVGD-EDA — bilan}
\author{}
\date{\today}

\begin{document}
\maketitle

\section{Objectif}

Tester si le critère inner-product test (Bollapragada, Byrd \& Nocedal, 2018) peut piloter
dynamiquement la taille de batch Monte-Carlo ($\lambda_a$) par agent SVGD-EDA, plutôt qu'un
$\lambda_a$ fixe.

Critère, recalculé à chaque tour de doubling sur les $\lambda_a$ échantillons déjà tirés :

\begin{itemize}
  \item $z^{(i)} = a^{(i)} \cdot \nabla_\theta \log \pi_\theta(x^{(i)})$ — contribution par
    échantillon ($a^{(i)}$ : avantage rang/fitness ; score en forme close $x-p$ ou
    $\mathrm{onehot}(x)-p$).
  \item $\hat g = \frac{1}{\lambda_a}\sum_i z^{(i)}$ — gradient estimé.
  \item $\widehat{\mathrm{tr}\Sigma} = \frac{1}{\lambda_a-1}\sum_i \|z^{(i)}-\hat g\|^2$ —
    dispersion (bruit Monte-Carlo).
  \item $\|\hat g\|^2_{\mathrm{corr}} = \max(\|\hat g\|^2 - \widehat{\mathrm{tr}\Sigma}/\lambda_a,\,0)$
    — norme débiaisée du bruit.
  \item $\hat g^\top \Sigma \hat g = \|\hat g\|^4 \cdot \mathrm{Var}(\hat s^{(i)})$, avec
    $\hat s^{(i)} = \langle z^{(i)}, \hat g\rangle/\|\hat g\|^2$ — variance projetée.
  \item $\lambda_{a,\mathrm{req}} = \hat g^\top \Sigma \hat g / (\mathrm{ip\_tol}^2 \cdot
    \|\hat g\|^4_{\mathrm{corr}})$ — taille requise ; $+\infty$ si $\|\hat g\|^2_{\mathrm{corr}}=0$.
\end{itemize}

Doubling ($\lambda_{a,\mathrm{init}} \to 2\times \to \dots \to \lambda_{a,\max}$) jusqu'à
$\lambda_a \geq \lambda_{a,\mathrm{req}}$ ou $\lambda_{a,\max}$. Si plafonné : shrinkage du pas
SVGD, $\eta \leftarrow \eta \cdot \|\hat g\|^2_{\mathrm{corr}} /
(\|\hat g\|^2_{\mathrm{corr}} + \widehat{\mathrm{tr}\Sigma}/\lambda_a)$.

\section{Implémentation}

\begin{itemize}
  \item Décision indépendante par (instance, agent) — pas d'agrégation sur le batch.
  \item $\lambda_{a,\mathrm{init}}=10$ par défaut (cf. §\ref{sec:lambdainit} pour la
    sensibilité), $\lambda_{a,\max}$ modeste (24-48).
  \item Contrat externe inchangé : \texttt{agent\_lambdas} reste constant à $\lambda_{a,\max}$,
    instances arrêtées tôt complétées par duplication de leurs échantillons réels.
\end{itemize}

\section{Résultats — NK (rugosité contrôlée par K)}

\subsection{Score final vs baseline $\lambda_a=10$ fixe}

Budget=50000, 10 instances $\times$ 10 restarts.

\begin{center}
\begin{tabular}{@{}lccc@{}}
\toprule
 & K=2 & K=4 & K=8 \\
\midrule
\textbf{baseline} & \textbf{0.7477} & \textbf{0.7608} & \textbf{0.7542} \\
meilleur adaptatif & 0.7466 (ip5/lmax24) & 0.7575 (ip5/lmax48) & 0.7408 (ip5/lmax48) \\
écart & $-0.15\%$ & $-0.4\%$ & $\mathbf{-1.8\%}$ \\
\bottomrule
\end{tabular}
\end{center}

L'écart au baseline grandit avec K, de façon monotone.

\subsection{Fraction plafonnée et shrinkage sur un run complet}

Run complet (50000 évaluations, $\lambda_{a,\mathrm{init}}=10$, $\lambda_{a,\max}=48$), mesuré
sur toutes les itérations cumulées. Colonnes : \%capped / shrink médian.

\begin{center}
\begin{tabular}{@{}llcccc@{}}
\toprule
dim & K & ip\_tol=0.4 & ip\_tol=5 & ip\_tol=20 & ip\_tol=100 \\
\midrule
64  & 1 & 89.0\% / 0.262   & 10.1\% / 0.0045 & 9.7\% / 0.0044  & 6.7\% / 0.0028 \\
64  & 2 & 92.5\% / 0.221   & 12.5\% / 0.0039 & 9.7\% / 0.0042  & 7.8\% / 0.0027 \\
64  & 4 & 98.9\% / 0.099   & 7.7\% / 0.0000  & 8.6\% / 0.0032  & 4.1\% / 0.0000 \\
64  & 8 & 99.8\% / 0.0013  & 14.2\% / 0.0000 & 8.9\% / 0.0000  & 7.8\% / 0.0000 \\
256 & 8 & 100.0\% / 0.0000 & 31.9\% / 0.0000 & 17.9\% / 0.0000 & 14.5\% / 0.0000 \\
\bottomrule
\end{tabular}
\end{center}

\begin{itemize}
  \item Rendements décroissants sur ip\_tol : chute nette entre 0.4 et 5, puis effet marginal
    ($5\to20\to100$), car ip\_tol n'intervient que si $\|\hat g\|^2_{\mathrm{corr}}>0$.
  \item Résidu incompressible qui empire avec la dimension : même à ip\_tol=100, 6-8\% restent
    plafonnées à dim=64 et \textbf{14.5\%} à dim=256 (K=8), avec shrink médian quasi nul.
\end{itemize}

\subsection{Effet de la dimension : NK N=256, K=8}

Même protocole, problème 4$\times$ plus grand (256 au lieu de 64) :

\begin{center}
\begin{tabular}{@{}lc@{}}
\toprule
config & score \\
\midrule
\textbf{baseline rang, $\lambda_a=10$} & \textbf{0.7235} \\
adaptatif ip5, lmax48  & 0.6162 \\
adaptatif ip20, lmax48 & 0.6401 \\
adaptatif ip50, lmax48 & 0.6423 \\
\textbf{écart (meilleur cas)} & \textbf{$-11.5\%$} \\
\bottomrule
\end{tabular}
\end{center}

Dégradation bien plus forte qu'à N=64 K=8 ($-1.8\%$) alors que le \%capped n'augmente que
modérément — la rugosité (K) et la dimension (N) semblent se cumuler plutôt que jouer
indépendamment.

\subsection{Sensibilité à $\lambda_{a,\mathrm{init}}$ (QUBO, rang, $\lambda_{a,\max}=48$)}
\label{sec:lambdainit}

\begin{center}
\begin{tabular}{@{}lcc@{}}
\toprule
$\lambda_{a,\mathrm{init}}$ & ip\_tol=5 & ip\_tol=20 \\
\midrule
\textbf{10} & $-301.60$ & $-295.00$ \\
3           & $-299.20$ & $-291.60$ \\
\bottomrule
\end{tabular}
\end{center}

$\lambda_{a,\mathrm{init}}=3$ fait systématiquement moins bien que $10$ : redescendre sous le
$\lambda_a$ connu-bon coûte plus qu'il ne rapporte en économie de budget.

\section{Résultats — QUBO (dim=64, instance t=5)}

Même reward des deux côtés (rang) — nb\_instances\_test=5, nb\_restarts=2, budget=50000 ;
\textbf{plus bas = meilleur} :

\begin{center}
\begin{tabular}{@{}lc@{}}
\toprule
config & score \\
\midrule
\textbf{baseline rang, $\lambda_a=10$} & \textbf{$-304.80$} \\
adaptatif ip5, lmax24  & $-304.20$ \\
adaptatif ip5, lmax48  & $-301.60$ \\
adaptatif ip20, lmax24 & $-298.80$ \\
adaptatif ip20, lmax48 & $-295.00$ \\
adaptatif ip50, lmax24 & $-304.00$ \\
adaptatif ip50, lmax48 & $-300.60$ \\
\bottomrule
\end{tabular}
\end{center}

Meilleur cas quasi identique au baseline (bruit d'un seul run) ; aucun cas ne le bat.

\textit{Note} : un premier essai avait suggéré un gain massif, dû à un facteur confondu (reward
différent entre adaptatif et baseline) ; une fois égalisé, le gain disparaît.

\section{Diagnostic}

\begin{enumerate}
  \item Paysage combinatoire intrinsèquement rugueux (surtout K élevé) : $\|\hat
    g\|^2_{\mathrm{corr}}$ reste souvent sous le bruit même avec beaucoup d'échantillons —
    régime courant, pas un cas rare de fin de convergence comme suppose le papier.
  \item Doubling multiplicatif : phénomène tout-ou-rien, $\lambda_a$ explose vite au plafond.
  \item Shrinkage quasi total sur les instances plafonnées, gelant une fraction croissante de
    la population à mesure que K augmente.
  \item SVGD-EDA tire sa force du nombre d'itérations (répulsion, moyennisation du bruit dans
    le temps), pas de la fiabilité de chaque pas — philosophie opposée à l'inner-product test.
\end{enumerate}

\section{Conclusion}

Une fois isolé de tout facteur confondu, le batch adaptatif \textbf{n'égale jamais un
$\lambda_a$ fixe bien réglé} : quasi neutre sur les paysages peu rugueux et petits (NK K=1-2,
QUBO), nettement dégradé sur les paysages rugueux et/ou grands (NK K=8, jusqu'à $-11.5\%$ à
N=256). La dégradation s'aggrave avec la rugosité (K) \textbf{et} la dimension (N), et ces deux
effets se cumulent. Le mécanisme fonctionne comme prévu par sa propre logique (plus de budget
là où le signal est faible), mais cette logique ne s'aligne pas avec les besoins d'un
optimiseur multi-agents itératif.

\end{document}
